\subsubsection{3.3.1 门店与订单对象访问矩阵}
\noindent 订单归属以服务器保存的原取货门店为准，登录角色和所属门店每次请求都重新核验。异店归还只是实物接收，不意味着外店人员取得顾客合同或财务资料。管理员的账号维护权亦不附带业务订单读取权。

\begin{BSRTable}
\rowcolor{TableHead}\color{white}\bfseries 执行者 & \color{white}\bfseries 订单资料、操作与例外 \\
顾客 & 只看本人预约、合同、归还鉴定和押金；他人订单及证据不可按编号查找。\\
门店租赁专员、门店运营经理 & 本店可看完整业务订单并在岗位允许时备货、核码交付、复核和结算。外店订单不进入列表、提醒、搜索及相关合同或照片接口。\\
异店接收人员 & 顾客出示完整订单号或取货码，且订单仍在租赁中时，仅看未归还单件的库存号、名称、尺码与归还截止；提交时再次核码，接收门店固定为当前员工所属门店。回执只含本批实收事实、时间及门店。\\
经营负责人、平台运维管理员 & 可看按门店和状态聚合的未解决/新增提醒及经营指标；不返回单笔订单、客户身份、取货码、押金、照片或合同。管理员仅另获账号和权限维护所需的元数据及权限审计。\\
装备检修技师 & 只看本人受派维修与鉴定任务；流程责任人中顾客身份匿名化，不获取订单号、支付或押金信息。\\
\end{BSRTable}

\BSRKeyPoint{授权边界}{按对象而非页面核验}{订单 ID、合同 ID、鉴定单与照片等直接或间接入口均由服务端重新判定；不存在与无权访问返回同类拒绝，不允许通过编号探测外店业务。}
\noindent 同一实物在异店归还时只改变当前所在门店；永久归属仍由经理发起调拨并由目标店接收。经理发布价目、采购入库、排班和调拨仅能作用于所属门店；跨店轮岗、临时代办和真实结算授权规则仍需 BSR 正式确认。
\noindent 门店人员名册仅列所属店员工与可受派技师候选；店员只读本人班次，经理只读所属店班次及每日排班，外店人员不能被排入本店。店员和经理的逐件库存限所属店；跨店供给协调读取按门店与状态聚合的数量。经营负责人和管理员仅得库存汇总，不取得外店库存号、单件成本或检查备注。
